\documentclass[11pt,a4paper]{article}
\usepackage[T1]{fontenc}
\usepackage[utf8]{inputenc}
\usepackage{lmodern}
\usepackage[margin=27mm]{geometry}
\usepackage{amsmath,amssymb,amsthm,mathtools}
\usepackage{microtype,booktabs,longtable,array,enumitem,xcolor}
\usepackage[colorlinks=true,linkcolor=blue!45!black,citecolor=blue!45!black,urlcolor=blue!50!black,breaklinks=true]{hyperref}
\usepackage{fancyhdr}
\pagestyle{fancy}\fancyhf{}\fancyhead[L]{\small Keio / Research manuscript}\fancyhead[R]{\small September 2026}\fancyfoot[C]{\thepage}
\setlength{\headheight}{14pt}
\setlength{\emergencystretch}{3em}
\setlength{\parskip}{0.35em}
\setlist{itemsep=2pt,topsep=4pt}
\newtheorem{theorem}{Theorem}[section]
\newtheorem{lemma}[theorem]{Lemma}
\newtheorem{proposition}[theorem]{Proposition}
\newtheorem{corollary}[theorem]{Corollary}
\newtheorem{conjecture}[theorem]{Conjecture}
\newtheorem{criterion}[theorem]{Criterion}
\theoremstyle{definition}\newtheorem{definition}[theorem]{Definition}
\newtheorem{example}[theorem]{Example}
\theoremstyle{remark}\newtheorem{remark}[theorem]{Remark}
\newtheorem{audit}[theorem]{Proof-status note}
\newcommand{\Q}{\mathbb Q}\newcommand{\Z}{\mathbb Z}\newcommand{\N}{\mathbb N}\newcommand{\R}{\mathbb R}\newcommand{\C}{\mathbb C}
\newcommand{\Qbar}{\overline{\mathbb Q}}
\DeclareMathOperator{\Tr}{Tr}\DeclareMathOperator{\Norm}{N}\DeclareMathOperator{\trdeg}{trdeg}\DeclareMathOperator{\ord}{ord}\DeclareMathOperator{\lcm}{lcm}
\DeclareMathOperator{\Ein}{Ein}\DeclareMathOperator{\Ei}{Ei}\DeclareMathOperator{\Ci}{Ci}\DeclareMathOperator{\Chi}{Chi}\DeclareMathOperator{\Li}{Li}
\newcommand{\floor}[1]{\left\lfloor #1\right\rfloor}\newcommand{\ceil}[1]{\left\lceil #1\right\rceil}
\newcommand{\sourceT}[1]{\textup{[source T#1]}}
\author{Keio}
\date{5 September 2026\\\small Standalone research draft; source audit identifiers retained}
\title{Dyadic Lucas Prime Cascades and\\ Composite Fibonacci Shifts}
\newcommand{\vTwo}{v_2}
\newcommand{\Fq}{\mathbb F}
\begin{document}
\maketitle
\begin{abstract}
Let $P$ be a positive odd integer, and let $U_0=0$, $U_1=1$,
$U_{n+2}=PU_{n+1}+U_n$. We prove that, for every integer $h$, the
set of $r\ge1$ for which $U_{2^r}+h$ is composite has positive lower
natural density. The proof gives a more precise propagation statement:
a sufficiently late prime value divides an arithmetic progression of
later values. An explicit threshold is obtained from the two quantities
$v_2((P^2+4)(h-1)^2+3)$ and $v_2((P^2+4)(h+1)^2+3)$.
The argument combines exact convergence of the even and odd dyadic
subsequences in $\Z_2$ with return of finite-field powers after their
2-primary orders have been exhausted. Taking $P=1$ answers Saito's
Problem~1.8 for every integer shift, with positive lower density in place
of infinitude. All sequence identities and valuation calculations used
in the proof are established here. The source identifiers T43 and T44
are retained only for cross-reference with the author's research notes;
historical priority beyond the cited literature is not claimed.
\end{abstract}

\section{The problem and the main theorem}
Saito's Problem~1.8 in the April 2025 version of
\cite{SaitoIntervals} asks whether, for each fixed integer $h$, the
Fibonacci values
\begin{equation}\label{eq:question}
 F_{2^r}+h\qquad(r\ge1)
\end{equation}
are composite infinitely often. Here $F_0=0$, $F_1=1$ and
$F_{n+2}=F_{n+1}+F_n$. The question keeps both the Fibonacci sequence
and the dyadic index $2^r$ fixed. We keep these quantifiers and answer
it through a general Lucas sequence result.

For $S\subseteq\N$, write
\[
 \underline d(S)=\liminf_{X\to\infty}
 \frac{\#(S\cap[1,X])}{\lfloor X\rfloor}.
\]
Throughout, a composite number is a positive integer with a proper
positive divisor other than $1$. Thus negative and zero values, or the
value $1$, are not called composite; each can occur only finitely often
in the increasing sequences considered below. For nonzero rational
numbers, $v_2$ denotes the exponent of $2$, normalized by $v_2(2)=1$.
It extends to $\Q_2$ in the usual way, with $v_2(0)=+\infty$.

Fix a positive odd integer $P$. Define
\begin{equation}\label{eq:Udef}
 U_0=0,\quad U_1=1,\quad U_{n+2}=PU_{n+1}+U_n,
 \qquad D=P^2+4.
\end{equation}
For an integer $h$, put
\begin{equation}\label{eq:Bdef}
 B_{P,h}=\max\left\{
 v_2\bigl(D(h-1)^2+3\bigr),
 v_2\bigl(D(h+1)^2+3\bigr)\right\}.
\end{equation}
The two arguments of $v_2$ are positive integers, so this threshold is
always finite and can be computed directly from $P$ and $h$.

\begin{theorem}[Prime cascade; \sourceT{43}]\label{thm:cascade}
Suppose that $r\ge B_{P,h}+1$ and
\[
 p=U_{2^r}+h
\]
is a prime satisfying $p\nmid2D$. There exists an integer $k\ge1$ such
that
\begin{equation}\label{eq:cascade}
 p\mid U_{2^{r+jk}}+h\qquad\text{for every }j\ge1.
\end{equation}
Every number in \eqref{eq:cascade} is larger than $p$ and is composite.
Consequently,
\begin{equation}\label{eq:density}
 \underline d\{r\ge1:U_{2^r}+h\text{ is composite}\}>0
\end{equation}
for every positive odd $P$ and every $h\in\Z$.
\end{theorem}

The step $k$ is effectively obtainable once a prime seed satisfying the
hypotheses is supplied: it is a multiplicative order in a finite cyclic
group. The density assertion does not assert that a prime seed exists.
If there are no sufficiently late prime values, all sufficiently late
values are already composite. This alternative is essential to the
proof and removes any need to prove the existence of primes in the
sequence.

\begin{corollary}[Fibonacci shifts; \sourceT{44}]\label{cor:Fib}
For every integer $h$,
\[
 \underline d\{r\ge1:F_{2^r}+h\text{ is composite}\}>0.
\]
In particular, the assertion posed in Saito's Problem~1.8 is true for
all integer shifts.
\end{corollary}

The main theorem has three ingredients. First, the dyadic Lucas values
have two opposite 2-adic limits. Second, those limits are roots of a
quadratic equation with no rational integer solutions; hence their
2-adic distances from the two relevant integer shifts have bounded
valuations. Third, these bounds make the remaining finite-field root
order odd. Repeated squaring then becomes invertible and returns to the
initial state. Sections~2--5 give the complete argument; Sections~6--8
record examples, certification procedures and the exact scope.

\section{The dyadic Lucas state}
Let $V_0=2$, $V_1=P$, and $V_{n+2}=PV_{n+1}+V_n$ be the companion
Lucas sequence. Set
\[
 f_r=U_{2^r},\qquad \ell_r=V_{2^r}\qquad(r\ge1).
\]
These two integer coordinates suffice for every part of the proof.

\begin{lemma}[Doubling identities]\label{lem:double}
For all $r\ge1$,
\begin{align}
 f_{r+1}&=f_r\ell_r, &
 \ell_{r+1}&=\ell_r^2-2,\label{eq:double}\\
 \ell_r^2-Df_r^2&=4.\label{eq:pell}
\end{align}
Moreover, $f_1=P$, $\ell_1=P^2+2$, every $f_r$ is odd, and
$\ell_r\equiv3\pmod4$.
\end{lemma}
\begin{proof}
Let $\alpha,\beta$ be the roots of $X^2-PX-1$. Then
$\alpha+\beta=P$, $\alpha\beta=-1$ and
$(\alpha-\beta)^2=D$. Direct substitution in the recurrence and its
initial conditions gives
\[
 U_n=\frac{\alpha^n-\beta^n}{\alpha-\beta},\qquad
 V_n=\alpha^n+\beta^n.
\]
The identities
\[
 U_{2n}=U_nV_n,\qquad
 V_{2n}=V_n^2-2(-1)^n,\qquad
 V_n^2-DU_n^2=4(-1)^n
\]
follow by expansion. Taking the even integer $n=2^r$ gives
\eqref{eq:double}--\eqref{eq:pell}. Since $P$ is odd,
$\ell_1=P^2+2\equiv3\pmod8$ and $f_1$ is odd.
If $\ell_r\equiv3\pmod4$, then $\ell_r^2-2\equiv3\pmod4$;
\eqref{eq:double} also preserves oddness of $f_r$.
\end{proof}

The real sizes cause no difficulty. The sequence $U_n$ is positive
for $n\ge1$ and strictly increasing for $n\ge2$. This follows directly
from its recurrence; the only initial equality is $U_1=U_2=1$ when
$P=1$. In particular, $f_r$ is strictly increasing for $r\ge1$ and
$f_r\to\infty$. These facts will turn divisibility into proper
divisibility after a prime seed.

\section{Two exact 2-adic limits}
The real growth of $f_r$ coexists with 2-adic convergence on its two
parity classes. We need the exact rates, not merely existence of the
limits.

\begin{lemma}[Valuation identities]\label{lem:valuations}
For every $r\ge1$,
\begin{align}
 v_2(\ell_r+1)&=r+1,\label{eq:ellval}\\
 v_2(f_{r+2}-f_r)&=r+1.\label{eq:fval}
\end{align}
\end{lemma}
\begin{proof}
Since $P^2\equiv1\pmod8$, one has
$v_2(\ell_1+1)=v_2(P^2+3)=2$. Also
\[
 \ell_{r+1}+1=(\ell_r-1)(\ell_r+1).
\]
The congruence $\ell_r\equiv3\pmod4$ gives
$v_2(\ell_r-1)=1$, proving \eqref{eq:ellval} by induction.
For the second identity, \eqref{eq:double} gives
\begin{align*}
 f_{r+2}-f_r
 &=f_r(\ell_r\ell_{r+1}-1)\\
 &=f_r(\ell_r^3-2\ell_r-1)\\
 &=f_r(\ell_r+1)(\ell_r^2-\ell_r-1).
\end{align*}
The first and last factors are odd. Equation~\eqref{eq:ellval}
therefore proves \eqref{eq:fval}.
\end{proof}

\begin{lemma}[Parity limits]\label{lem:limits}
There exist odd elements $u_0,u_1\in\Z_2$, with $u_1=-u_0$, such that
$f_r\to u_i$ through $r\equiv i\pmod2$. They satisfy
\begin{equation}\label{eq:uroot}
 Du_i^2=-3
\end{equation}
and the exact approximation formula
\begin{equation}\label{eq:exactrate}
 v_2(f_r-u_{r\bmod2})=r+1.
\end{equation}
Neither $u_0$ nor $u_1$ is a rational integer.
\end{lemma}
\begin{proof}
For either fixed parity, \eqref{eq:fval} shows that the successive
differences have valuations tending to infinity. The corresponding
subsequence is Cauchy in the complete ring $\Z_2$. Its limit is odd
because all terms are odd. In the convergent series
\[
 u_{r\bmod2}-f_r
 =\sum_{j\ge0}(f_{r+2j+2}-f_{r+2j}),
\]
the first summand has valuation $r+1$ and all remaining summands have
valuation at least $r+3$. The ultrametric inequality, with a unique
term of least valuation, gives \eqref{eq:exactrate}.
Equation~\eqref{eq:ellval} says $\ell_r\to-1$. Passing to the limit
in $f_{r+1}=f_r\ell_r$ along either parity gives $u_1=-u_0$.
Passing to the limit in \eqref{eq:pell} gives $1-Du_i^2=4$, which is
\eqref{eq:uroot}. If $u_i=c\in\Z$, this equality would hold in
$\Q$ and hence in $\R$, contradicting $D>0$ and $Dc^2\ge0$.
\end{proof}

No choice of a real square root of $-3/D$ is involved. The roots in
Lemma~\ref{lem:limits} are constructed in $\Z_2$ by the recurrence.
Their failure to be ordinary integers is precisely what provides a
uniform distance bound for each fixed shift.

\begin{lemma}[Distance from an integer]\label{lem:distance}
Let $u\in\{u_0,u_1\}$ and $c\in\Z$. Then
\begin{equation}\label{eq:distance}
 v_2(u-c)+v_2(u+c)=v_2(Dc^2+3).
\end{equation}
In particular, both valuations on the left are finite, nonnegative,
and at most the quantity on the right.
\end{lemma}
\begin{proof}
The product identity
\[
 (u-c)(u+c)=u^2-c^2=-\frac{Dc^2+3}{D}
\]
and the oddness of $D$ give \eqref{eq:distance}. Neither factor is zero
by Lemma~\ref{lem:limits}; both belong to $\Z_2$.
\end{proof}

\section{From a shifted prime to an odd root order}
The relevant finite-field root order divides either $p-1$ or
$2(p+1)$. The next proposition explains why their 2-parts become
bounded when $p$ itself is a sufficiently late shifted sequence value.

\begin{proposition}[The order budget]\label{prop:budget}
Let $r\ge B_{P,h}+1$ and suppose $p=f_r+h$ is an odd prime. Then
\begin{equation}\label{eq:budget}
 v_2(p-1)\le B_{P,h},\qquad v_2(p+1)\le B_{P,h}.
\end{equation}
\end{proposition}
\begin{proof}
Put $u=u_{r\bmod2}$, $c_-=1-h$ and $c_+=-1-h$. Then
$p-1=f_r-c_-$ and $p+1=f_r-c_+$. Lemma~\ref{lem:distance} gives
\[
 v_2(u-c_\pm)\le B_{P,h}.
\]
On the other hand, \eqref{eq:exactrate} gives
$v_2(f_r-u)=r+1>B_{P,h}$. Thus, in
$f_r-c_\pm=(f_r-u)+(u-c_\pm)$, the second summand has strictly
smaller valuation. The valuation of the sum equals that of this
summand, proving \eqref{eq:budget}.
\end{proof}

\begin{lemma}[Finite-field root orders]\label{lem:orders}
Let $p\nmid2D$ be prime. A root $\gamma$ of $X^2-PX-1$ in
$\Fq_{p^2}$ has multiplicative order $t$ satisfying
\[
 \begin{cases}
 t\mid p-1,&D\text{ is a square modulo }p,\\
 t\mid2(p+1),&D\text{ is a nonsquare modulo }p.
 \end{cases}
\]
If $p=f_r+h$ and $r\ge B_{P,h}+1$, the order of
$z=\gamma^{2^r}$ is odd.
\end{lemma}
\begin{proof}
The constant term is nonzero modulo $p$, so $\gamma\ne0$.
If $D$ is a square modulo $p$, the roots lie in $\Fq_p^\times$,
whose order is $p-1$. If $D$ is a nonsquare, the Frobenius automorphism
interchanges the two distinct roots. Since their product is $-1$,
\[
 \gamma^p=-\gamma^{-1},\qquad \gamma^{p+1}=-1.
\]
Thus $t\mid2(p+1)$. Proposition~\ref{prop:budget} now implies
$v_2(t)\le B_{P,h}+1\le r$. The order of $\gamma^{2^r}$ equals
$t/\gcd(t,2^r)$ and has no factor $2$.
\end{proof}

\section{Return and positive lower density}
We can now complete both the propagation statement and its density
consequence without using any distribution theorem for prime values
of Lucas sequences.

\begin{proof}[Proof of Theorem~\ref{thm:cascade}]
Assume first the prime-seed hypotheses. Choose a root $\gamma$ as in
Lemma~\ref{lem:orders}, put $z=\gamma^{2^r}$, and let $q$ be its odd
multiplicative order. If $q=1$, set $k=1$; otherwise set
$k=\ord_q(2)$. In either case, for every $j\ge1$,
\[
 z^{2^{jk}}=z,
 \qquad\text{hence}\qquad
 \gamma^{2^{r+jk}}=\gamma^{2^r}.
\]
The other root is $\delta=-\gamma^{-1}$. All exponents under
consideration are even, so the same equality holds with $\delta$ in
place of $\gamma$. The finite-field version of Binet's formula is
valid because $p\nmid D$:
\[
 U_n=\frac{\gamma^n-\delta^n}{\gamma-\delta}
 \quad\text{in }\Fq_{p^2}.
\]
It follows that $U_{2^{r+jk}}\equiv U_{2^r}\pmod p$, proving
\eqref{eq:cascade}. Strict increase gives
$U_{2^{r+jk}}+h>U_{2^r}+h=p$, so the divisor is proper.

For the density conclusion, choose $R$ so large that all $f_r+h$,
$r\ge R$, exceed $1$, exceed every prime divisor of $2D$, and satisfy
$r\ge B_{P,h}+1$. This is possible because $f_r\to\infty$.
If no $f_r+h$ with $r\ge R$ is prime, every such value is composite
and the composite-index set has density $1$. Otherwise choose one
prime seed $p=f_r+h$, $r\ge R$. The progression
$r+k,r+2k,\ldots$ from \eqref{eq:cascade} lies in the composite-index
set and has natural density $1/k$. Therefore the latter set has lower
natural density at least $1/k>0$.
\end{proof}

\begin{proof}[Proof of Corollary~\ref{cor:Fib}]
The recurrence \eqref{eq:Udef} with $P=1$ is exactly the Fibonacci
recurrence. Thus Theorem~\ref{thm:cascade} applies with $D=5$ and
\[
 B_{1,h}=\max\{v_2(5(h-1)^2+3),v_2(5(h+1)^2+3)\}.
\]
It treats every fixed $h\in\Z$ and leaves the exponent $2^r$
unchanged. Positive lower density implies infinitude.
\end{proof}

\begin{remark}[The parity shortcut]
When $h$ is odd, every $f_r+h$ is even and eventually greater than
$2$, so the composite-index set is cofinite. The finite-field argument
is needed for even shifts. Keeping both parities in the theorem makes
the quantifier over all integers $h$ explicit.
\end{remark}

\section{Worked exact examples}
\subsection{A seed beyond the explicit threshold}
Take $P=1$, $h=-16$. Then
\[
 5(h-1)^2+3=1448=8\cdot181,
 \qquad5(h+1)^2+3=1128=8\cdot141,
\]
so $B_{1,-16}=3$. At $r=4$,
\[
 F_{16}-16=987-16=971.
\]
The integer $971$ is prime: trial division by the primes at most
$\sqrt{971}<32$ suffices. Modulo $971$, the roots of
$X^2-X-1$ are $174$ and $798$. The first has order $485$, and
\[
 174^{16}\equiv481\pmod{971},\qquad
 \ord_{971}(481)=485,\qquad\ord_{485}(2)=48.
\]
These equalities can be checked by modular powering, using
$485=5\cdot97$ and $48=2^4\cdot3$ to test the proper prime-factor
quotients of each claimed order. Therefore
\begin{equation}\label{eq:971}
 971\mid F_{2^{4+48j}}-16\qquad(j\ge1).
\end{equation}
This is an explicit infinite arithmetic progression of composite
indices and gives the concrete lower bound $1/48$ for this shift.
The bound need not be optimal.

\subsection{A short return below the sufficient threshold}
For $h=2$, one has $B_{1,2}=4$. The values
\[
 F_8+2=23,\qquad F_{16}+2=989=23\cdot43
\]
exhibit divisibility before the sufficient threshold $r\ge5$.
This does not contradict the theorem; the threshold is sufficient,
not necessary. There is also a direct infinite certificate. At $r=4$,
the dyadic state modulo $23$ is
\[
 (f_4,\ell_4)\equiv(21,-1)\pmod{23}.
\]
The second coordinate remains $-1$ under
$\ell\mapsto\ell^2-2$, while the first changes sign at each step.
Hence $f_r\equiv21\pmod{23}$ for every even $r\ge4$, and
\[
 23\mid F_{2^r}+2\qquad(r\ge4\text{ even}).
\]
The composite-index set for $h=2$ therefore has lower density at least
$1/2$ by this separate finite-state certificate.

\subsection{Why the late-seed hypothesis cannot simply be dropped}
Take $P=1$, $h=4$, $r=2$. The value $F_4+4=7$ is prime, but it never
divides a later shifted dyadic Fibonacci value. Indeed,
\[
 (f_2,\ell_2)=(3,7)\equiv(3,0)\pmod7.
\]
The next first coordinate is zero and stays zero under
$f\mapsto f\ell$. Thus
\[
 F_{2^s}+4\equiv4\pmod7\qquad(s\ge3).
\]
This disproves the unqualified assertion that every prime seed
propagates. Here $B_{1,4}=7$, so $r=2$ is well outside the theorem's
hypotheses. A proof that suppresses the order budget would miss this
example.

\section{A finite certificate for a supplied seed}
Once $P,h,r,p$ are given with $p=U_{2^r}+h$, the infinite divisibility
statement admits a finite mathematical certificate. The following
steps describe that certificate, not a search for arbitrary late prime
values.
\begin{enumerate}
\item Certify that $p$ is prime and that $p\nmid2(P^2+4)$.
\item Evaluate \eqref{eq:Bdef} and verify $r\ge B_{P,h}+1$.
\item Construct $\Fq_p$ or $\Fq_{p^2}$ according as the discriminant
is split or nonsplit, and choose a root $\gamma$ of $X^2-PX-1$.
\item Compute $z=\gamma^{2^r}$, its odd order $q$, and a positive
integer $k$ with $2^k\equiv1\pmod q$; the exact order of $2$ is
optional if a larger valid $k$ is easier to certify.
\item The identities in Section~5 certify every future index
$r+jk$, so no further integer factorization is needed.
\end{enumerate}
The root computation can also be replaced by a quotient-ring
calculation modulo $p$. In either representation, a certificate is
verified by finitely many exact arithmetic operations. A successful
finite certificate proves a whole progression; checking a large number
of isolated terms would not provide the same conclusion.

\begin{proposition}[A convenient universal return step]
Under the prime-seed hypotheses, let
\[
 W=\frac{2(p^2-1)}{2^{v_2(2(p^2-1))}}.
\]
If $W=1$ set $k=1$; otherwise any positive multiple of
$\varphi(W)$ may be used in \eqref{eq:cascade}.
\end{proposition}
\begin{proof}
The root order divides $2(p^2-1)$ in either splitting case. After
raising to $2^r$, its remaining odd order $q$ divides $W$.
Since $W$ is odd, Euler's theorem gives
$2^{\varphi(W)}\equiv1\pmod W$, hence modulo $q$. The proof of
Theorem~\ref{thm:cascade} now applies.
\end{proof}
The step supplied by this proposition can be much larger than the
least return step. Its purpose is to show that the certificate does
not require a minimal-order calculation.

\section{Scope, limitations and provenance}
The proved assertion is about composite shifted values, not about the
frequency of prime values. It allows the possibility that only
finitely many shifted values are prime, and it does not prove that
the composite-index set has a natural density. Lower density is the
precise conclusion supported by the argument.

The restriction to odd $P$ supplies the odd dyadic state and the exact
valuation law \eqref{eq:ellval}. No assertion for all even $P$ is made
by this proof. Likewise, the recurrence has second characteristic
coefficient $-1$ in the conventional Lucas notation $U_n(P,-1)$.
Changing that coefficient changes the Pell identity and the limiting
equation, so a separate argument is required.

The proof does not settle the infinitude of composite Fermat numbers.
Its decisive input is a shift-dependent bound for the 2-primary
parts of $p-1$ and $p+1$. For Fermat-type candidates those parts may
grow with the exponent, and the odd-order conclusion need not hold.
Nor does the Fibonacci specialization settle Saito's distinct cubic
Pisot problem involving $\floor{\beta^{3^r}}$, or determine the
transcendence of Mills' constant.

The general theorem corresponds to T43 in the source research notes;
Corollary~\ref{cor:Fib} corresponds to T44. The doubling formulas,
finite-field root orders and valuation rules are classical algebraic
tools, proved explicitly above to make this manuscript independent.
The contribution presented here is their assembled prime-seed
propagation statement and its application to the exact published
question. This draft does not identify a theorem as historically first
merely because a literature search did not locate an identical result.
Independent expert review and a fuller priority search remain
appropriate before a priority claim is attached to publication.

\begin{thebibliography}{9}
\bibitem{SaitoIntervals}
K.~Saito,
\emph{Intervals without primes near an iterated linear recurrence
sequence}, arXiv:2504.14968, version~1, 21 April 2025,
Problem~1.8.
\href{https://arxiv.org/abs/2504.14968}{arXiv:2504.14968}.
Published in \emph{International Journal of Number Theory} (2026),
\href{https://doi.org/10.1142/S1793042126500685}
{doi:10.1142/S1793042126500685}.
\end{thebibliography}
\end{document}
